% ==========================================================================
%  Chapter 50 — Phase 5: Dynamic \& Seismic Analysis Verification
% ==========================================================================

\chapter{Phase 5: Dynamic and Seismic Analysis Verification}
\label{ch:phase5-dynamic-seismic}

This chapter documents Phase~5 of the master plan
(Chapter~\ref{ch:convergence-master-plan}).
The original scope required implementing the full time-stepping
pipeline: mass matrix assembly, Rayleigh damping, PETSc~TS
integration with the generalized-$\alpha$ scheme, ground-motion
input, and per-step material state commitment.
A thorough audit revealed that the \texttt{DynamicAnalysis} class
was \emph{already fully implemented} (approximately 770~lines in
\texttt{src/analysis/DynamicAnalysis.hh}), together with the
boundary-condition infrastructure in
\texttt{src/model/BoundaryCondition.hh} and the damping factories
in \texttt{src/analysis/Damping.hh}.

Phase~5 was therefore focused on a comprehensive \emph{verification
campaign} exercising nine distinct aspects of the pipeline, bringing
the number of passing dynamic-analysis assertions from 62 (existing
\texttt{test\_dynamic\_analysis.cpp}) to $62+36=98$.


%% ------------------------------------------------------------------
\section{Phase~5 Scope and Audit}
\label{sec:phase5-scope}

Chapter~\ref{ch:convergence-master-plan} listed six deliverables for
Phase~5:
\begin{enumerate}
  \item Mass matrix assembly (consistent, lumped, density API).
  \item Rayleigh damping: $\mathbf{C}=\alpha_M\mathbf{M}+\beta_K\mathbf{K}$.
  \item PETSc~TS generalized-$\alpha$ time stepper.
  \item Ground-motion pseudo-force: $-\mathbf{M}\hat{\mathbf{g}}\,a_g(t)$.
  \item Per-step monitor callback for state commitment.
  \item VTK time-series output via~\texttt{PVDWriter}.
\end{enumerate}

All six were found operational.  The implementation uses PETSc's
\texttt{I2Function}/\texttt{I2Jacobian} interface, which formulates
the second-order ODE residual
\begin{equation}
  \mathbf{F}(t,\mathbf{u},\dot{\mathbf{u}},\ddot{\mathbf{u}})
  = \mathbf{M}\ddot{\mathbf{u}}
  + \mathbf{C}\dot{\mathbf{u}}
  + \mathbf{f}_\mathrm{int}(\mathbf{u})
  - \mathbf{f}_\mathrm{ext}(t)
  = \mathbf{0},
\end{equation}
and the system Jacobian
\begin{equation}
  \mathbf{J} = a\,\mathbf{M} + v\,\mathbf{C} + \mathbf{K}_t(\mathbf{u}),
\end{equation}
where $a$ and~$v$ are shift parameters supplied by PETSc's
\texttt{TSALPHA2} integrator.


%% ------------------------------------------------------------------
\section{DynamicAnalysis Architecture}
\label{sec:phase5-architecture}

\Cref{tab:phase5-architecture} summarises the key classes
and their roles.

\begin{table}[htbp]
\centering
\caption{Dynamic-analysis component inventory.}
\label{tab:phase5-architecture}
\begin{tabular}{lp{7.5cm}}
\toprule
Component & Role \\
\midrule
\texttt{DynamicAnalysis<Policy>}
  & Owns PETSc \texttt{TS}, mass matrix~$\mathbf{M}$,
    damping matrix~$\mathbf{C}$, and Jacobian work
    matrix~$\mathbf{J}$.  Registers \texttt{I2Function}
    and \texttt{I2Jacobian} callbacks, monitors, and
    ground-motion pseudo-forces. \\
\texttt{BoundaryConditionSet<dim>}
  & Aggregates initial conditions, nodal forces
    (time-dependent), prescribed displacements,
    and ground-motion records. \\
\texttt{GroundMotionBC}
  & Uniform base acceleration in one direction via
    $-\mathbf{M}\hat{\mathbf{g}}_\mathrm{dir}\,a_g(t)$. \\
\texttt{NodalForceBC<dim>}
  & Time-dependent nodal force with arbitrary
    \texttt{TimeFunction} per component. \\
\texttt{damping::*}
  & Factory functions: \texttt{rayleigh},
    \texttt{rayleigh\_from\_frequencies},
    \texttt{rayleigh\_from\_single\_ratio},
    \texttt{mass\_proportional},
    \texttt{stiffness\_proportional},
    \texttt{none}. \\
\texttt{PVDWriter}
  & VTK Parallel-VTU Data (.pvd) time-series writer. \\
\bottomrule
\end{tabular}
\end{table}

The generalized-$\alpha$ method (\texttt{TSALPHA2}) offers a single
user-facing parameter: the spectral radius at infinity
$\rho_\infty \in [0,1]$.  Setting $\rho_\infty=1$ yields the
trapezoidal rule (no numerical dissipation); $\rho_\infty=0$
maximises high-frequency damping while preserving second-order
accuracy.


%% ------------------------------------------------------------------
\section{Verification Campaign}
\label{sec:phase5-tests}

The verification tests reside in
\texttt{tests/test\_phase5\_dynamic\_verification.cpp} and use a
single-element Hex8 unit cube ($8$~nodes, $24$~DOFs) with
$E=1000$, $\nu=0$, $\rho=1$, and $x=0$ face clamped.
\Cref{tab:phase5-test-results} presents the nine test groups.

\begin{table}[htbp]
\centering
\caption{Phase~5 verification results (36~assertions, all passing).}
\label{tab:phase5-test-results}
\begin{tabular}{clcc}
\toprule
\# & Test description & Assertions & Status \\
\midrule
1 & Free vibration period (zero-crossing + CSV)   & 4 & \checkmark \\
2 & Undamped amplitude conservation                & 4 & \checkmark \\
3 & Rayleigh damping amplitude decay               & 4 & \checkmark \\
4 & Damping factory completeness (6~factories)     & 6 & \checkmark \\
5 & Ground-motion seismic input                    & 4 & \checkmark \\
6 & Multi-directional ground motion ($x+z$)        & 3 & \checkmark \\
7 & Forced vibration with CSV output               & 4 & \checkmark \\
8 & J$_2$ plasticity under dynamic loading         & 3 & \checkmark \\
9 & Spectral-radius comparison ($\rho_\infty=1$ vs~$0$) & 4 & \checkmark \\
\midrule
  & \textbf{Total}                                 & \textbf{36} & \\
\bottomrule
\end{tabular}
\end{table}


\subsection{Test~1: Free Vibration Period}
\label{sec:phase5-t1}

An initial axial displacement $u_0=10^{-3}$ is applied at the free
face.  The monitor callback records the time history of $u_x$ at
node~1, and the period is computed from zero-crossings.  For the
1D~bar approximation:
\begin{equation}
  \omega_1 = \frac{\pi}{2L}\sqrt{\frac{E}{\rho}}
  \approx 49.67\ \mathrm{rad/s},
  \qquad
  T_1 \approx 0.127\ \mathrm{s}.
\end{equation}
The measured period from the 3D~Hex8 element is $T\approx 0.115$~s,
which is lower due to the three-dimensional stiffness contribution
but physically consistent.  A CSV file is written for visual
inspection.


\subsection{Test~2: Undamped Amplitude Conservation}
\label{sec:phase5-t2}

Without Rayleigh damping ($\alpha_M=\beta_K=0$), the generalized-%
$\alpha$ method with $\rho_\infty=1$ is non-dissipative.  The
displacement envelope remained within $[{-1.000}\times 10^{-3},\,
{+1.000}\times 10^{-3}]$ after 150~steps, confirming energy
conservation with numerical tolerance $<0.1\%$.


\subsection{Test~3: Rayleigh Damping Decay}
\label{sec:phase5-t3}

Strong Rayleigh damping ($\alpha_M=10$, $\beta_K=0.002$) produces a
decay ratio (late-quarter peak / early-quarter peak) of~$\approx
0.10$, demonstrating $\sim 90\%$ amplitude reduction over three
analytical periods.  The time history is written to CSV for
inspection.


\subsection{Test~4: Damping Factories}
\label{sec:phase5-t4}

All five \texttt{damping::} factory functions return non-null
callables, and \texttt{damping::none()} returns an empty callable
(evaluates to \texttt{false}).  This validates the full Rayleigh
damping API surface.


\subsection{Test~5: Ground Motion Seismic Input}
\label{sec:phase5-t5}

A sinusoidal ground-acceleration pulse
$a_g(t) = 0.5\sin(30t)$ is applied in the $x$-direction for
$t \le 0.2$~s, followed by free vibration.  The maximum relative
displacement exceeds $3\times 10^{-4}$, confirming that the
pseudo-force vector $-\mathbf{M}\hat{\mathbf{g}}\,a_g(t)$ is
correctly assembled and applied.


\subsection{Test~6: Multi-Directional Ground Motion}
\label{sec:phase5-t6}

Two simultaneous ground motions ($x$~and~$z$) with distinct
amplitudes ($A_x=0.3$, $A_z=0.5$) and frequencies
($\omega_x=20$, $\omega_z=35$~rad/s) are applied.  Both
displacement components develop non-trivial response, verifying
that the stacking of multiple \texttt{GroundMotionBC} objects
in the \texttt{BoundaryConditionSet} is functional.


\subsection{Test~7: Forced Vibration with CSV Output}
\label{sec:phase5-t7}

A harmonic nodal force $F(t) = F_0 \sin(\omega_f t)$ at
sub-resonant frequency ($\omega_f = \omega_1/2$) is distributed
across the free-face nodes.  The test verifies convergence,
non-trivial forced response, and that the CSV file is written
with $\ge 20$~timestep records.


\subsection{Test~8: J$_2$ Plasticity Dynamic}
\label{sec:phase5-t8}

A soft J$_2$ material ($E=200$, $\sigma_y=0.05$, $H=5$) is loaded
with a large initial displacement $u_0=0.01$ to guarantee yielding.
The monitor callback invocation count verifies per-step material
state commitment, and the final displacement remains bounded,
confirming that the TS solver handles nonlinear plasticity under
dynamic loading without divergence.


\subsection{Test~9: Spectral Radius Comparison}
\label{sec:phase5-t9}

Two otherwise-identical simulations differ only in spectral radius:
$\rho_\infty=1.0$ (no dissipation) and $\rho_\infty=0.0$ (maximum
high-frequency dissipation).  Both converge successfully and produce
non-trivial amplitudes, validating the
\texttt{TSAlpha2SetRadius} control path.


%% ------------------------------------------------------------------
\section{Full Regression Status}
\label{sec:phase5-regression}

After Phase~5, the complete \texttt{ctest} suite shows:

\begin{center}
\begin{tabular}{lr}
\toprule
Metric & Value \\
\midrule
Total tests     & 46 \\
Passing         & 44 \\
Pre-existing failures & 2 \\
\quad\texttt{cantilever\_benchmark} & (numerical tolerance) \\
\quad\texttt{snes\_gmsh\_pipeline}  & (missing mesh file) \\
New regressions & 0 \\
\bottomrule
\end{tabular}
\end{center}

The Phase~5 test (\texttt{phase5\_dynamic\_verification}) completes
all 36~assertions in approximately 18~seconds.


%% ------------------------------------------------------------------
\section{Summary}
\label{sec:phase5-summary}

Phase~5 verified the dynamic-analysis pipeline that was already
fully implemented in \texttt{fall\_n}:

\begin{itemize}
  \item \textbf{770-line} \texttt{DynamicAnalysis} header with
    PETSc~TS generalized-$\alpha$, I2Function/I2Jacobian callbacks,
    Rayleigh damping, ground-motion pseudo-forces, and per-step
    monitor/commit.
  \item \textbf{BoundaryConditionSet} infrastructure with composable
    \texttt{TimeFunction} factories (constant, harmonic, ramp,
    step, pulse, piecewise-linear, exponential decay, product,
    sum, scale).
  \item \textbf{6~damping factories} covering Rayleigh,
    frequency-based, single-ratio, mass-proportional,
    stiffness-proportional, and null.
  \item \textbf{98~total dynamic-analysis assertions} across two
    test files, covering free vibration, amplitude conservation,
    damping decay, ground motion (single and multi-directional),
    forced vibration with CSV output, nonlinear J$_2$ plasticity,
    and spectral-radius control.
\end{itemize}

Phases~0--5 are now complete.  The remaining deliverable is
Phase~6: MPI scalability audit.
